\documentclass[10pt,letterpaper]{article} % exam class defaults to 10pt on letterpaper
\usepackage[margin=1in]{geometry}         % exam class defaults to 1-inch margins
\usepackage{graphicx} 
\usepackage{fancyhdr}
\usepackage{amssymb}
\usepackage{amsmath}
\usepackage{mathtools}
\usepackage[most]{tcolorbox}
\usepackage[dvipsnames]{xcolor}
\usepackage{blindtext}
\usepackage{amsthm}
\usepackage{upgreek}
\usepackage{hyperref}

\setlength{\parindent}{0pt}
\setlength{\parskip}{1ex plus 0.5ex minus 0.2ex} % Matches the paragraph spacing of exam

\newtheorem{thm}{Theorem}[section]
\newtheorem{lem}{Lemma}[section]
\newtheorem{defn}{Definition}[section]
\newtheorem{cor}{Corollary}[section]
\newtheorem{axm}{Axiom}[section]

\tcolorboxenvironment{thm}{colback=lime, colframe=black, breakable}
\tcolorboxenvironment{lem}{colback=yellow, colframe=black, breakable}
\tcolorboxenvironment{defn}{colback=pink, colframe=black, breakable}
\tcolorboxenvironment{cor}{colback=SpringGreen, colframe=black, breakable}
\tcolorboxenvironment{axm}{colback=white, colframe=black, breakable}

% --- Header/Footer Setup ---
\fancyhf{}
\renewcommand\headrulewidth{0.5pt}
\pagestyle{fancy}
\author{Omkar Tasgaonkar}
\setlength{\headsep}{0.5in}
\fancyfoot[C]{(\thepage)}
\renewcommand\footrulewidth{0.5pt}

\fancypagestyle{firstpage}{
    \fancyhf{}
    \fancyhead[C]{\textbf{Notes on Analysis}}
    \renewcommand\headrulewidth{0.5pt}
    \fancyfoot[C]{(\thepage)}
    \renewcommand\footrulewidth{0.5pt}
}
\begin{document}
\tableofcontents
\newpage
\thispagestyle{firstpage}

\section{Characterizations of Limits and Continuity}
The characterization of limits and continuity we will do will branch off the two following topological definitions:
\begin{enumerate}
\item \textbf{Adherent Point}: $\forall r > 0: B(x, r) \cap A \neq \varnothing$
\item \textbf{Limit Point}: $\forall r > 0: B(x, r)\setminus \{x\} \cap A \neq \varnothing$
\end{enumerate}
We start off by formally defining a limit as the following. $L$ is a limit for a function $f$ at $x \in \overline{\text{Dom}(f)}$ if:
\begin{equation*}
  \forall \epsilon > 0, \exists \delta > 0, \forall y \in X: 0 < \varrho(y, x) < \delta \implies \varrho(f(y), L) < \epsilon
\end{equation*}
Notice two things: first is that the existence of a limit does not care about what the value of $f(x)$ is (if it is defined at all). The second thing is we require $x$ to be in the closure of the domain, not the domain
\section{Density}
We start off by considering some topological space $X$, endowed with a topology we will call $\uptau_X$. Then for some $A \subset X$, we have the following notion of density:
\begin{defn}
A subset of a topological space $A \subseteq X$ is dense in $X$ if (we propose two equivalent conditions):
\begin{enumerate}
  \item $\forall x \in X, \forall U_x \in \uptau_X: U \cap A \neq \varnothing$
  \item $\bar{A} = X$
\end{enumerate}
\end{defn}
Intuitively, no matter how close we look to a point in $X$, we can find a point in $A$, so $A$ "covers" (perhaps an interesting adjective and intuively useful is "spans") all of $X$. We call this variety of dense as \textbf{everywhere dense}.\\

It however is possible that the closure of $A$ is not ALL of $X$, but rather a subset of $X$. We will call this dense in a subset:
\begin{defn}
Consider $A \subseteq Y \subseteq X$, then $A$ is dense in $Y$ if $Y \subseteq \bar{A}$.
\end{defn}
Note that we did not give the open set definition; for the purposes of further content dealing with closures is easier. Examining the condition $Y \subseteq \bar{A}$, we see that implies by definition of closure that $\bar{Y} \subseteq \bar{A}$. Furthermore, since $A \subseteq Y$, we have that $\bar{A} = \bar{Y}$. This is an equivalent condition, which will be indispensable moving forward, and we will default to this. Additionally, using the open set condition for dense in a subset requires an additional notion:
\begin{defn}
The subspace topology is an induced topology on a subset $Y \subseteq X$, where:
\begin{equation*}
  \uptau_Y \coloneq \{O \cap Y : O \in \uptau_X\}
\end{equation*}
\end{defn}
An open set relative to $Y$ (we call this relatively open as well), is open if and only if it can be represented as the intersection of an open set $O \subseteq X$ and $Y$. We talked about $A$ being dense in a subset meaning the closure of $A$ is the closure of the subset. The closure we referred to was the closure of the sets with respect to the entire space $X$.\\

Consider the set $\mathbb{Q}$ in $\mathbb{Q}$ versus $\mathbb{Q}$ in $\mathbb{R}$. Then the closure of $\mathbb{Q}$ in the first case is $\mathbb{Q}$ itself; whereas, the closure in $\mathbb{R}$, is $\mathbb{R}$. Then:
\begin{lem}
The closure of $A$ with respect to $Y$, which we will denote as $\overline{A}^Y$ is the same as $\bar{A} \cap Y$.
\end{lem}
\begin{proof}
We know that the definition of closure with respect to $X$ is:
\begin{equation*}
  \bar{A} \coloneq \bigcap \{C: C \ \text{closed} \land A \subseteq C\}
\end{equation*}
Then for all such $C$ which are closed subsets of $Y$, there exists some $B \subseteq X$ such that $B$ is closed and $C = B \cap Y$. Then:
\begin{equation*}
  \bar{A}^Y = \bigcap C = \bigcap (B \cap Y) = \left(\bigcap B\right) \cap Y = \bar{A} \cap Y
\end{equation*}
\end{proof}
\begin{lem}
For $A$ dense in a subset $Y$, another equivalent form of the statement is $\bar{A}^Y = Y$.
\end{lem}
\begin{proof}
We first prove that $\bar{A} = \bar{Y}$ yields the equivalent statement. Simply take the intersection with $Y$ for both sets such that:
\begin{equation*}
  \bar{A} \cap Y = \bar{A}^Y = \bar{Y} \cap Y = Y
\end{equation*}
Now if we have that $\bar{A}^Y = Y$, we can write it as $\bar{A} \cap Y = Y$. Thus:
\begin{equation*}
  Y = \bar{A} \cap Y \subseteq \bar{A} \implies \bar{Y} \subseteq \bar{A}
\end{equation*}
By subset properties once more, we get $\bar{A} = \bar{Y}$.
\end{proof}
Notice how in Definition 0.2, we required $A \subseteq Y$. We can actually weaken this condition to simply a nonempty intersection:
\begin{defn}
$A$ is \textbf{somewhere dense} if there exists some open $Y \subseteq X$ such that $\overline{A \cap Y} = \bar{Y}$
\end{defn}
Consider some open set $Y$ such that $O \coloneq A \cap Y$ and $A \cap Y \neq \varnothing$. Then $O \subseteq Y$, so if $O$ is dense in $Y$:
\begin{equation*}
  \bar{O}^Y = Y \iff \bar{O} = \bar{Y} \iff \overline{A \cap Y} = \bar{Y}
\end{equation*}
This property of somewhere dense means part of $A$ can span some subset, not necessarily the whole set. We give a few facts without apology:
\begin{enumerate}
  \item A subset $A$ is everywhere dense (we will refer to this as simply dense) if and only if somewhere dense for all open subsets $Y$ of $X$
  \item A subset $A$ is somwhere dense if and only if $\text{int}(\bar{A}) \neq \varnothing$
\end{enumerate}
We introduce the notion, jumping off of somewhere dense, of nowhere dense.
\begin{defn}
A subset $A$ is nowhere dense if and only if it is NOT somewhere dense.
\end{defn}
Notice by the (1) in the previous list, not dense \textbf{DOES NOT} imply nowhere dense. If a set is not dense, then there exists some $Y$ in which it is not somewhere dense (but the lack of somewhere density is not necessarily true for all such open subsets).\\

However, a nowhere dense subset is not dense. By (2) we see that a subset $A$ is nowhere dense if and only if $\text{int}(\bar{A}) = \varnothing$. This will be the main criterion for proving nowhere density as we will see with examples such as the Cantor set. We state an additional corollary:
\begin{cor}
If $A$ is somewhere dense for some open $Y \subseteq X$, then $\overline{A \cap Y} = \bar{A} \cap \bar{Y}$
\end{cor}
We already have the inclusion that $\overline{A \cap Y} \subseteq \bar{A} \cap \bar{Y}$ (which is true for all sets). We prove the reverse inclusion. We see that:
\begin{equation*}
  \bar{Y} = \overline{A \cap Y} \land A \cap Y \subseteq A \implies \bar{Y} \subseteq \bar{A} 
\end{equation*}
Thus $\bar{A} \cap \bar{Y} = \bar{Y} = \overline{A \cap Y}$, actually directly proving equality.

\section{The Cantor Set}

\subsection{Properties}
The Cantor Set is a nowhere dense, perfect set constructed by repeatedly removed the middle $1/3$ interval of each interval. We denote each step of this iterative process as $C_n$ where for all $I \subset C_{n-1}$, we remove their middle 1/3 portion, creating $C_n$.\\

We first note that the Cantor set is nonempty. Since for all $n$, we have that $C_{n+1} \subseteq C_n$ and because $[0, 1]$ is closed and bounded (compact), by Cantor intersection property we have that their intersection, which is $C$, is nonempty.\\

An interesting property of the Cantor Set is that it is uncountable. For this, we examine a representation of the elements of the Cantor Set, called the ternary expansion:
\begin{equation*}
  \forall x \in C: x \coloneq \sum_{i \in \mathbb{N}} \frac{2\sigma_i}{3^{i+1}}
\end{equation*}
where $\sigma_x$ is a sequence in $\{0, 1\}^\mathbb{N}$. This is called the ternary representation because it is in base 3. If a representation of a number in $[0, 1]$ MUST contain a $1$ in it, then that tells us it is in a middle 1/3 interval.\\

Based on its place in the decimal, such a ternary expansion tells us that in order to reach the number, we must pass through the interval $[1/3^n, (2/3)^n)$ (and endpoints remain in the Cantor Set), so it belongs in the middle interval.\\

This ternary expansion gives us an injection $f: \{0, 1\}^\mathbb{N} \xrightarrow{} C$. If $f(\sigma_1) = f(\sigma_2)$, then:
\begin{equation*}
  f(\sigma_1) = f(\sigma_2) \implies \sum_{i \in \mathbb{N}}\frac{2\sigma_{1i}}{3^{i+1}} = \sum_{i \in \mathbb{N}}\frac{2\sigma_{2i}}{3^{i+1}}
\end{equation*}
\begin{equation*}
  \implies \sum_{i \in \mathbb{N}}\frac{2(\sigma_{1i} - \sigma_{2i})}{3^{i+1}} = 0
\end{equation*}
Since $0$ has a unique ternary expansion which is $0.00000..._3$ itself, then this means $\forall i \in \mathbb{N}: \sigma_{1i} = \sigma_{2i}$, proving equality of the sequences and injectivity of $f$.\\

Now because $\{0, 1\}^\mathbb{N} \simeq \mathbb{R}$, then $\mathbb{R} \lesssim C$, so the Cantor set is uncountable (which will be an interesting result later on).\\

We note that each time the iterative process doubles the amount of intervals, so each $C_n$ can be covered by $2^n$ closed balls of radius $3^{-n}/2$. This fact will be quite important when we examine integration over the indicator function on $C$. We first prove that the Cantor set is nowhere dense. 
\begin{lem}
The Cantor Set is nowhere dense.
\end{lem}
First since each $C_n$ is a closed set, then the construction of $C_{n+1}$ involves removing only OPEN middle 1/3 sets, so $C_{n+1}$ is also a finite union of $2^n$ closed sets, thus closed. $C$ is also closed because it is the arbitrary union of closed sets $C_n$, so $\bar{C} = C$.\\

A point is in the interior if of a set $A$:
\begin{equation*}
  \exists \epsilon > 0: B(x, \epsilon) \subseteq A
\end{equation*} 
So we aim to prove the negation of this for all such points in the Cantor set. For all $x$ in the Cantor set, since $x \in \cap_{n \in \mathbb{N}} C_n$, then $\forall n \in \mathbb{N}: x \in C_n$. We can find the following:
\begin{equation*}
  \forall \epsilon > 0, \exists n \in \mathbb{N}: \frac{1}{3^n} < 2\epsilon
\end{equation*}
Then we have that $x$ belongs to some interval $I_{n'}$ in $C_n$. Because each interval is covered by a closed ball of radius $1/(2 \cdot 3^n)$ centered at the midpoint (which may be $x$ or not), then the ball $B(x, \epsilon)$ certainly exceeds the interval.\\

(If $x$ is the midpoint, this follows from the covering. If $x$ is not the midpoint, then the ball $B(x, 1/(2 \cdot 3^n))$ already falls outside the interval, so we can extend it to the $\epsilon$-ball by subset properties). Thus:
\begin{equation*}
  \forall \epsilon > 0, \exists y \notin C: y \in B(x, \epsilon)
\end{equation*}
Thus the Cantor set has no interior points and is nowhere dense.
\begin{lem}
The Cantor Set is perfect, which means it is closed and has no isolated points.
\end{lem}
Given some $x \in C$ and some $\epsilon > 0$, the previous lemma gave us the existence of $n \in \mathbb{N}$, for which $x$ belongs to an interval in $C_n$. Once again by case work, we can show that at least one endpoint of the interval different from $x$ should fall in the ball $B(x, \epsilon)$ and since endpoints remain in $C$, we have shown that no point is isolated:
\begin{equation*}
  \forall x \in C, \forall \epsilon > 0: (B(x, \epsilon)\setminus \{x\}) \cap C \neq \varnothing
\end{equation*}

\smallskip
\subsection{Integrability}
We now examine the following function:
\begin{equation*}
  1_C \coloneq \begin{cases*}
    1 \ : \ x \in C\\
    0 \ : \ x \notin C
  \end{cases*}
\end{equation*}
We aim to show this function is Riemann integrable. We first show that it is discontinuous at all points within the Cantor set (via limsup liminf criterion):
\begin{equation*}
  \forall x \in C, \forall \delta > 0: \sup_{z: |z - x| < \delta}f(x) = 1 \land \inf_{z: |z - x| < \delta}f(x) = 0
\end{equation*}
The infimum part is true because of Lemma 0.3 (nowhere dense), which tells us that for any $\delta$-ball around $x$, we can always find some point outside the Cantor set within the ball. The supremum can also be supported by the perfect + nowhere dense properties of $C$ (if we consider the punctured open ball, perhaps to disprove existence of a limit). In that case, we have a discontinuity of second kind at each $x \in C$ as we can show both left and right limits do not exist.\\

Then since the above condition is true for all $\delta$:
\begin{equation*}
  \lim_{\delta \xrightarrow{} 0+}\sup_{z: z \in B(x, \delta)}f(x) = 1 > \lim_{\delta \xrightarrow{} 0^+}\inf_{z: |z-x| < \delta}f(x) = 0
\end{equation*}
which tells us that $f$ is not continuous at all $x \in C$. So now we have a function that is nowhere dense, with an uncountable number of discontinuities, that too of the second kind YET is Riemann integrable. Then:
\begin{enumerate}
  \item Cardinality of the set of discontinuities while sufficient is too strong of a condition for Riemann integrability
  \begin{enumerate}
    \item Thomae's function in $[0, 1]$ showed that we can have countably many discontinuities (but of the first kind as limit exists at all rationals)
  \end{enumerate}
  \item Type of discontinuities faces the same issue as well ($1_C$ has discontinuities of the second kind)
\end{enumerate}

\begin{lem}
The function $1_C$ is bounded and Riemann integrable on $[0, 1]$ and $\int_{0}^{1}1_C dx = 0$.
\end{lem}
We note that by the Darboux definition, the lower integral over any partition is always $0$. So we aim to show that the upper integral also converges to $0$. We define the following:
\begin{equation*}
  \Pi_n \coloneq \{I_n \cup M_n: I_n \subseteq C_n\}
\end{equation*}
where $I_n$ are the existing intervals after removing the middle 1/3 intervals denoted as $M_n$. Then:
\begin{equation*}
  \forall \Pi_n: U(f, \Pi_n) = \sum_{I_{n_i} \in I_n}1 \cdot \Delta I_{n_i} + \sum_{M_{n_j} \in M_n} 0 \cdot \Delta M_{n_j}
\end{equation*} 
This is because the supremum on $I_n$ is always 1 (the endpoints are in $C$) and the supremum on $M_n$ is 0 (it contains no points in the Cantor set). Then:
\begin{equation*}
  U(f, \Pi_n) = 1 \cdot \left(\frac{2}{3}\right)^n \implies \inf U(f, \Pi_n) = \lim_{n \xrightarrow{} 0}U(f, \Pi_n) = 0
\end{equation*}
Thus the upper integral and lower integral agree at the value $0$. The $(2/3)^n$ come from the fact that there are $2^n$ closed balls of radius $1/(2 \cdot 3^n)$ that cover all of $C_n$. Notice that the "length" of the Cantor set (the set of discontinuities in this case) goes to 0.

\medskip
\section{Measure}

\subsection{Definitions}
We never officially introduced the notion of length. We first say that the length of any OPEN interval $(a, b)$, (denoted by $l$) in $\mathbb{R}$ is $b-a$. Length we will consider to be a map $l: \mathcal{P}(\mathbb{R}) \xrightarrow{} [0, \infty]$ (range in extended reals). We now introduce length of a set as the following:
\begin{defn}
The outer measure of a set $A$ is via covering sequences of intervals $\{I_n\}_{n \in \mathbb{N}}$:
\begin{equation*}
  |A| \coloneq \inf\left\{\sum_{n \in \mathbb{N}}l(I_n) : A \subseteq \bigcup_{n \in \mathbb{N}} I_n\right\}
\end{equation*}
\end{defn}
We note that the measure of the empty set is $0$. The intuition behind this is consider all possible covers for $A$, then the smallest possible cover (the one that "hugs" or matches $A$'s shape the "tightest"), can best approximate its length. The idea of measure will give us the equivalent criterion for integrability, which we will cover later on.\\

Another way we can restate this property is as the following:
\begin{equation*}
  \forall \epsilon > 0, \exists \{I_n\}_{n \in \mathbb{N}}: \sum_{n \in \mathbb{N}}I_n - |A| < \epsilon
\end{equation*} 
First we examine some properties of measure:

\begin{lem}
The measure of a finite set of points is $0$.
\end{lem}
\begin{proof}
Let our finite set of points be $A \coloneq \{x_0, ..., x_m\}$. Then we can cover $A$ by the following sequence:
\begin{equation*}
  I_n \coloneq \begin{cases*}
    (x_n - \epsilon, x_n + \epsilon) \ : \ n \leqslant m\\
    \varnothing \ : \ n > m
  \end{cases*}
\end{equation*}
\begin{equation*}
  |A| \leqslant \sum_{n \in \mathbb{N}}l(I_n) = 2m\epsilon
\end{equation*}
Because the length of an interval is always non-negative, we note that the measure also must be non-negative. If we assume some set were to have measure less than $0$, then by properties of infimum being adherent, there would exist a sequence of covering sets such that their length is negative (not possible). Since $\epsilon$ is an arbitrary number greater than $0$, we can find any value of $\epsilon$ less than a possible value of $|A|$, thus squeezing the measure so that $|A| = 0$. 
\end{proof}
This proof and non-negativity of the metric tells us that it too is a map from $\mathcal{P}(\mathbb{R}) \xrightarrow{} [0, \infty]$. We present a stronger version of Lemma 0.6:
\begin{lem}
The measure of a countable set of points is 0.
\end{lem}
\begin{proof}
Consider a similar setup as Lemma 0.6, with $A \coloneq \{x_0, x_1,...\}$ and a specific $\{I_n\}_{n \in \mathbb{N}}$ defined as:
\begin{equation*}
  I_n \coloneq (x_n - 2^{-n-1}, x_n + 2^{-n-1})
\end{equation*}
Then:
\begin{equation*}
  |A| \leqslant \sum_{n \in \mathbb{N}}I_n = \sum_{n \in \mathbb{N}}\frac{\epsilon}{2^n} = 2\epsilon
\end{equation*}
The result comes from the sum of the geometric series of $(1/2)^n$. Once again we have a similar squeeze argument, showing $|A| = 0$. This is of interest to us as it shows that Thomae's function in $[0, 1]$ has a set of discontinuities of zero measure.
\end{proof}

So far we have only talked about measure in terms of open intervals $I_n$, however we now show that it is equivalent for closed intervals. We prove the following two Lemmas:
\begin{lem}
If $A \subseteq B$, then the measure of $A$ is less than or equal to the measure of $B$.
\end{lem}
\begin{proof}
Consider some sequence $\{I_n\}_{n \in \mathbb{N}}$ such that the union of the intervals cover $B$. Then:
\begin{equation*}
  A \subseteq \bigcup_{n \in \mathbb{N}}I_n \implies |A| \leqslant \sum_{n \in \mathbb{N}}l(I_n)
\end{equation*}
By Definition 0.6 (infimum properties), we get that $|A|$ is a lower bound for all such sums, thus $|A| \leqslant |B|$.
\end{proof}
\begin{lem}
The measure of an open interval $(a, b)$ is equal to the measure of the closed interval $[a, b]$.
\end{lem}
\begin{proof}
We defined the measure of the open interval $(a, b)$ as its length $b-a$. Then consider:
\begin{equation*}
  \forall \epsilon > 0: [a, b] \subseteq (a - \epsilon, b + \epsilon)
\end{equation*}
By Lemma 0.8, we see that:
\begin{equation*}
  |[a, b]| \leqslant |(a- \epsilon, b+\epsilon)| = (b-a) + 2\epsilon
\end{equation*}
Since $\epsilon$ is arbitrary, this tells us that $|[a, b]| \leqslant b-a = |(a, b)|$. Lemma 0.8 also yields the converse:
\begin{equation*}
  (a, b) \subseteq [a, b] \implies |(a, b)| = b-a \leqslant |[a, b]|
\end{equation*}
Thus $|[a, b]| = |(a, b)| = b-a$. This allows us to define measure in terms of closed intervals $\bar{I_n}$ as well as they are of equivalent length.
\end{proof}

Lemma 0.8 allows us to further argue that the measure of the set of discontinuities (or the measure of the Cantor set) for the function in Lemma 0.5 is 0:
\begin{equation*}
  \forall n \in \mathbb{N}: C \subseteq C_n \implies |C| \leqslant |C_n|
\end{equation*}
We know that $2^n$ closed balls of radius $1/(2\cdot 3^{n})$ cover the set so that $C_n \ \cup I_n$. The measure of each $C_n$ can be calculated by the measure of each interval $I_n$, for which we require the lemma:
\begin{lem}
Outer measure is subadditive so that:
\begin{equation*}
  \forall n \in \mathbb{N}: \left|\bigcup_{n \in \mathbb{N}}A_n\right| \leqslant \sum_{n \in \mathbb{N}}|A_n|
\end{equation*}
\end{lem}
This is a triangle inequality-"ish" property, while not exactly as at one side we have set unions and at the other side we have addition. Let $\textit{I}_n$ be a sequence of intervals covering its respective $A_n$, and $I_{ni}$ denoting an element in the sequence covering $A_n$. Then:
\begin{equation*}
  \bigcup_{n \in \mathbb{N}} A_n \subseteq \bigcup_{n \in \mathbb{N}}\textit{I}_n = \bigcup_{n \in \mathbb{N}}\bigcup_{m \in \mathbb{N}}I_{nm}
\end{equation*}
Since each $\textit{I}_n$ is a cover for $A_n$, we have that:
\begin{equation*}
  \forall n \in \mathbb{N}: |A_n| \leqslant \sum_{m \in \mathbb{N}}l(I_{nm}) \implies \sum_{n \in \mathbb{N}} |A_n| \leqslant \sum_{n \in \mathbb{N}}\sum_{m \in \mathbb{N}}l(I_{nm})
\end{equation*}
Now for all $\epsilon > 0$, by the equivalent definition of measure (derived from supremum), we can find sequences $\textit{I}_n$, such that:
\begin{equation*}
  \sum_{m \in \mathbb{N}}l(I_{nm}) \leqslant \sum_{n \in \mathbb{N}}|A_n| + \frac{\epsilon}{2^n}
\end{equation*} 
\begin{equation*}
  \implies \sum_{n \in \mathbb{N}}\sum_{m \in \mathbb{N}}l(I_{nm}) \leqslant \sum_{n \in \mathbb{N}}|A_n| + 2\epsilon
\end{equation*}
We now need to show the existence of a covering sequence for $\cup A_n$. Consider the sequence where the indices $nm$ add up to a specific number (starting from $0$), so that we get:
\begin{equation*}
  \{I_{00}\}, \{I_{01}, I_{10}\}, \{I_{11}, I_{20}, I_{02}\}...
\end{equation*}
Then $\forall \epsilon > 0$, there exists a sequence as defined above such that:
\begin{equation*}
    \left|\bigcup_{n \in \mathbb{N}}A_n\right| \leqslant \sum_{n \in \mathbb{N}}\sum_{m \in \mathbb{N}}l(I_{nm}) \leqslant \sum_{n \in \mathbb{N}}|A_n| + 2\epsilon
\end{equation*}
Since it is true for all epsilon, then we arrive at out result. This however only proves subadditivity (which in our case to prove the zero measure of the Cantor set would be enough, but we can prove something stronger).\\
\begin{lem}
For a sequence of disjoint open intervals $I_n$, we have that:
\begin{equation*}
  \left|\bigcup_{n \in \mathbb{N}}I_n\right| = \sum_{n \in \mathbb{N}}|I_n|
\end{equation*}
\end{lem}
$I_n$ is covering sequence for itself (and also we can infer by Lemma 0.10) the inequality:
\begin{equation*}
  \left|\bigcup_{n \in \mathbb{N}}I_n\right| \leqslant \sum_{n \in \mathbb{N}}|I_n|
\end{equation*}
By the definition of infimum and measure, we also have that:
\begin{equation*}
  \forall \epsilon > 0, \exists \{I_n'\}_{n \in \mathbb{N}}: \sum_{n \in \mathbb{N}}|I_n'| \leqslant \left|\bigcup_{n \in \mathbb{N}}I_n\right| + \epsilon
\end{equation*}
In order to prove this, we need to prove disjoint additivity in the finite sense, such that:
\begin{equation*}
  \left|\bigcup_{n = 0}^{N}I_n\right| = \sum_{n = 0}^{N}l(I_n)
\end{equation*}
Once again by definition of measure, the forward inequality holds. We prove the reverse inequality by induction. Our base case is when $N = 0$, in which it follows from the definition of length and measure. Now for $N+1$, we once again have by definition:
\begin{equation*}
  \left|\bigcup_{n = 0}^{N+1}I_n\right| \leqslant \sum_{n = 0}^{N+1}l(I_n)
\end{equation*}
On the other hand:
\begin{equation*}
  \sum_{n = 0}^{N+1}l(I_n) = \sum_{n = 0}^{N}l(I_n) + l(I_{n+1}) = \left|\bigcup_{n=0}^{N}l(I_n)\right| + l(I_{n+1})
\end{equation*}
Now let $I_n'$ be an arbitrary covering sequence for the union of $I_n$ up to $N+1$, the define:
\begin{align*}
  J_n &\coloneq I_n' \cap I_{N+1}\\
  J_n' &\coloneq I_n' \cap (\mathbb{R} \setminus I_{N+1})
\end{align*}
Because we have the condition of disjointness, we are able to separate the cover $I_n'$ into these two sequences such that $I_n' = J_n \cup J_n'$. We note that $J_n$ exclusively covers $I_{N+1}$ and does not intersect with any of the other sets, while $J_n'$ exclusively covers the rest of the union. Therefore:
\begin{equation*}
  \sum_{n \in \mathbb{N}} |I_n'| = \sum_{n \in \mathbb{N}}(|J_n| + |J_n'|) \geqslant \left|\bigcup_{n = 0}^{N}I_n\right| + |I_{N+1}|
\end{equation*}
Once again since it is a covering sequence for the union up to $N+1$, for any such sequence, the right hand side is a lower bound, so by infimum properties, we have that:
\begin{equation*}
  \left|\bigcup_{n = 0}^{N+1}I_n\right| \geqslant \left|\bigcup_{n= 0}^{N}I_n\right| + |I_{N+1}|
\end{equation*}
thereby proving equality.\\

Now returning to the Cantor set, by Lemma 0.11, we see that since each $C_n$ is a closed union of disjoint intervals of length $3^{-n}$:
\begin{equation*}
  \forall n \in \mathbb{N}: |C_n| = \left(\frac{2}{3}\right)^n \implies \forall n \in \mathbb{N}: |C| \leqslant \left(\frac{2}{3}\right)^n
\end{equation*}
By a similar argument to Lemma 0.7, we see that $n \xrightarrow{} 0 \implies |C| = 0$. If we say $|C| = \epsilon > 0$, then $\exists n \in \mathbb{N}: (2/3)^n < \epsilon$, which would mean $|C_n| < |C|$ (contradiction). Now the reader may think that the Cantor set is nowhere dense, so nowhere dense sets have zero measure; however, that is false. First consider the converse:
\begin{equation*}
  \text{zero measure} \nRightarrow \text{nowhere dense}
\end{equation*}
For this consider Thomae's function in the interval $[0, 1]$. The discontinuties (of the first kind) are at $\mathbb{Q} \cap [0, 1]$. Since this is a countable set, we have that its measure is $0$. However:
\begin{equation*}
  \overline{\mathbb{Q} \cap [0, 1]} = \bar{\mathbb{Q}} \cap [0, 1] = [0, 1]
\end{equation*}
So the set of discontinuities is dense in all of $[0, 1]$ (this also shows rationals are somewhere dense within the interval by Defn 0.4). For the converse, we introduce a more intricate idea of the Fat Cantor set.

\subsection{The Fat Cantor Set}
The Fat Cantor Set or rather the family of such sets are examples of sets which are nowhere dense but DO NOT have measure of 0. We consider a sequence $\{a_n\}_{n \in \mathbb{N}} \in (0, 1)^\mathbb{N}$, such that:
\begin{equation*}
  \sum_{n \in \mathbb{N}}a_n < \infty
\end{equation*}
We can now calculate the measure of each stage of the Fat Cantor set:
\begin{equation*}
  |F_n| \coloneq \prod_{n \in \mathbb{N}}(1 - \alpha_n)
\end{equation*}
To prove that this product does not diverge or converge to $0$, we can rewrite it in terms of sums as the following:
\begin{equation*}
  \prod_{n \in \mathbb{N}}(1 - a_n) = e^{\ln\left(\prod_{n \in \mathbb{N}}(1-a_n)\right)} = e^{\sum_{n \in \mathbb{N}}\ln(1-a_n)}
\end{equation*}
Now the convergence of the product then depends on the convergence of the sum of $\ln$ of its terms. If the sum converges to some value $L$, then the product converges to sum non-zero positive value (it only converges to zero if the sum diverges to $-\infty$). Now observe the power series for $\ln(1-x)$ around $x = 0$:
\begin{equation*}
  \ln(1-x) < P_2(x) \coloneq 0 - x - \frac{x^2}{2}
\end{equation*}
The less than inequality comes from the concavity of $\ln$ (any tangent line/graph falls above it), so we see that $\ln(1-x) < -x$ for all $x \in (0, 1)$. Thus:
\begin{equation*}
  \sum_{n \in \mathbb{N}}\ln(1-a_n) < \sum_{n \in \mathbb{N}} - a_n
\end{equation*}
and because we are given that the sum of $a_n's$ converge, then the sum on the right hand side also converges. Thus as $n \xrightarrow{} \infty, |F_n| \xrightarrow{} L > 0$, so the Fat Cantor set's measure is NONZERO. We now show that the Fat Cantor Set is nowhere dense. We see that the length of each interval $I_m^{(n)}$ in the nth stage can be bounded above:
\begin{equation*}
  |I_m^{(n)}| = \frac{1}{2^n}\prod_{n \in \mathbb{N}}(1- a_n) \leqslant \frac{1}{2^n}
\end{equation*}
So given some point $x \in F$, then given some $\epsilon > 0$, we can always find $n \in \mathbb{N}$ such that $2^{-n} < \epsilon$, which means:
\begin{equation*}
  \exists y \in B(x, \epsilon): y \notin F_n \implies y \notin F
\end{equation*}
So the Fat Cantor set is an example of a set that is nowhere dense yet not zero measure. Thus the two notions are separate (we showed the set of discontinuities in Thomae's function are dense but zero measure).

\pagebreak
\section{Fundamental Theorem of Calculus}
We first begin by defining the notion of an antiderivative as the following:
\begin{defn}
Given a function $f: [a, b] \xrightarrow{} \mathbb{R}$, we define the antiderivative $F: [a, b] \xrightarrow{} \mathbb{R}$ as a function that is continuous on $[a, b]$ and differentiable on $(a, b)$, such that:
\begin{equation*}
  \forall x \in (a, b): F'(x) = f(x)
\end{equation*}
\end{defn}
Note that the definition does not tell us that $f$ is Riemann integrable. Now the continuity of $F$ will play a key role in proving the second Fundamental Theorem of Calculus.\\

First, we see the following lemma:
\begin{lem}
If $f$ is Riemann Integrable on $[a, b]$, then $F(x) \coloneq \int_{a}^{x} f(t) dt$ and $F(a) \coloneq 0$ is lipschitz continuous on $(a, b)$ and left/right continuous at $a, b$ respectively.
\end{lem}
\begin{proof}
If $f$ is Riemann Integrable, then it is bounded on $[a, b]$. So:
\begin{equation*}
  \forall x, y \in (a, b): |F(x) - F(y)| = \left|\int_{y}^{x}f(t)dt\right| \leqslant ||f|| \cdot |x - y|
\end{equation*}
\end{proof}

Now this still does not tell us that our definition of $F(x)$ serves as a valid antiderivative. We further propose the following lemma, which provides us a sufficient condition to do so:
\begin{thm}
FTC1: If $f$ is continuous on $[a, b]$, then $\forall x \in (a, b): F'(x) = f(x)$.
\end{thm}
\begin{proof}
We see the following holds:
\begin{equation*}
  F(y) - F(x) = \int_{x}^{y}f(t)dx \implies F(y) - F(x) - f(x)(y-x) = \int_{x}^{y}(f(t)-f(x))dt
\end{equation*}
\begin{equation*}
  |y-x|\cdot \left|\frac{F(y) - F(x)}{y-x} - f(x)\right| = |y-x| \cdot \left|\int_{x}^{y}(f(t)-f(x))dt\right| \leqslant ||f(t) - f(x)||
\end{equation*}
\end{proof}
Since $f$ is continuous, we note that setting $y = x + \delta$,  $\inf_{\delta > 0} \sup_{t \in B(x, \delta)}|f(t) - f(x)| = 0$, showing that $F'(x) = f(x)$. We now prove the second Fundamental Theorem of Calculus:
\begin{thm}
FTC2: Given an antiderivative $F: [a, b] \xrightarrow{} \mathbb{R}$ of the function $f$ with $F'$ being Riemann Integrable, then:
\begin{equation*}
  \int_{a}^{b} F'(x)dx = F(b) - F(a)
\end{equation*}
\end{thm}
\begin{proof}
Since $F'$ is Riemann Integrable, we have that, considering the integral evaluates to the value $L$:
\begin{equation*}
  \forall \epsilon > 0, \exists \delta > 0, \forall \Pi: ||\Pi|| < \delta \implies |R(f, \Pi) - L| < \epsilon
\end{equation*}
If we can show that there exists a partition $\Pi$ with mesh less than $\delta$ and Riemann Sum $F(b) - F(a)$, then we have proven the theorem.\\

We choose the partition $\Pi \coloneq \{[a + i \frac{b-a}{n}, a + (i+1)\frac{b-a}{n}]: 0 \leqslant i \leqslant n-1\}$, where $\delta \cdot n > b-a$. 
\begin{equation*}
  F(b) - F(a) = \sum_{i=1}^{n}(F(t_i) - F(t_{i-1}))
\end{equation*}
Since $F$ is continuous on $[a, b]$ and differentiable on $(a, b)$, then, we can repeatedly apply MVT, so that:
\begin{equation*}
  = \sum_{i=1}^{n}(f(t_i^*)(t_i - t_{i-1}))
\end{equation*}
Then because $F'$ is Riemann Integrable, we get that:
\begin{equation*}
  \forall \epsilon: \left|F(b) - F(a) - \int_{a}^{b} F'(x)dx\right| < \epsilon
\end{equation*}
\end{proof}
Note that we require initial function to be Riemann Integrable. The existence of an antiderivative does not immediately give us this. For example, with Volterra's function, its derivative is discontinuous at the points in the Fat Cantor set.\\

Since the set has non-zero measure, then by Lebesgue Vitali criterion, the derivative is not integrable, so neither FTC1 nor FTC2 hold. 

\begin{cor}
U-Substitution: Given some continuous function $f$ on $[c, d]$ and a continuous function $\phi: [a, b] \xrightarrow{} (c, d)$ which is continuous on $[a, b]$ and differentiable on $(a, b)$:
\begin{equation*}
  \int_{\phi(b)}^{\phi(a)} f(t) dt = \int_{a}^{b}f(\phi(x)) \cdot \phi'(x) dx
\end{equation*}
\end{cor}
\begin{proof}
  
\end{proof}

\subsection{Variation}
We define variation over an interval of $f$ as:
\begin{equation*}
  V(f, [a, b]) \coloneq \sup_{\Pi}\sum_{i=1}^{n}|f(t_i) - f(t_{i-1})|
\end{equation*}
Now considering $f$ to be continuous on $[a, b]$ and $F$ to be the antiderivative of $f$, we want to show:
\begin{equation*}
  V(F, [a, b]) = \int_{a}^{b}|f(x)|dx
\end{equation*}
We prove the forward inequality first, which is the simpler direction:
\begin{equation*}
  \forall \Pi: \sum_{i=1}^{n}|F(t_i) - F(t_{i-1})| = \sum_{i=1}^n \left|\int_{t_{i-1}}^{t_i} f(x) dx \right|
\end{equation*}
In order to prove the inequality, we must use the following lemma:
\begin{lem}
The integral inequality states for all Riemann integrable functions $f$ on the interval $[a, b]$:
\begin{equation*}
  \left|\int_{a}^{b} f(x)dx \right| \leqslant \int_{a}^{b} |f(x)| dx
\end{equation*}
\end{lem}
\begin{proof}
We see that:
\begin{equation*}
  \forall \Pi_{[a, b]}: |R(f, \Pi)| = \left|\sum_{i=1}^{n} f(t_i^*)(t_i - t_{i-1})\right| \leqslant \sum_{i=1}^{n}|f(t_i^*)|(t_i - t_{i-1}) = R(|f|, \Pi)
\end{equation*}
Since it holds true for all Riemann Sums, taking limits, we have proved the integral inequality.
\end{proof}
Now returning to the original problem, we get that:
\begin{equation*}
  \sum_{i=1}^n \left|\int_{t_{i-1}}^{t_i} f(x) dx \right| \leqslant \sum_{i=1}^n \int_{t_{i-1}}^{t_i}|f(x)|dx = \int_{a}^{b}|f(x)|dx
\end{equation*}
Since this holds true for all partitions, we have found an upper bound, thus proving that $V(f, [a, b])$ is less than or equal to the integral of $|f|$ by supremum properties. We now prove the opposite direction:\\

Because $f$ is continuous and the composition of continuous functions is continuous, then $|f|$ is continuous thus Riemann Integrable. So:
\begin{equation*}
  \forall \epsilon > 0, \exists \delta > 0, \forall \Pi: ||\Pi|| < \delta \implies \left|\int_{a}^{b}|f(x)|dx - R(|f|, \Pi)\right| < \epsilon
\end{equation*}
Now, for each $\delta$ we take a partition where $n\delta > b-a$ defined as $\Pi \coloneq \{[a+i \cdot \frac{b-a}{n}, a + (i+1)\cdot \frac{b-a}{n}] : 0 \leqslant i \leqslant n-1\}$. Then applying MVT to each interval, we get a tagged partition $\Pi^*$:
\begin{equation*}
  \sum_{i=1}^{n}|f(t_i^*)|(t_i - t_{i-1}) = \sum_{i=1}^{n}|F(t_i) - F(t_{i-1})|
\end{equation*}
Thus by integrability, we have that:
\begin{equation*}
  \forall \epsilon > 0, \exists \Pi: \left|\int_{a}^{b}|f(x)|dx - \sum_{i=1}^{n}|F(t_i) - F(t_{i-1})|\right| < \epsilon
\end{equation*}
\begin{equation*}
  \implies \forall \epsilon > 0, \exists \Pi: \int_{a}^{b}|f(x)|dx \leqslant \sum_{i=1}^n|F(t_i) - F(t_{i-1})| + \epsilon
\end{equation*}
By supremum properties of variation
\begin{equation*}
  \implies \forall \epsilon > 0: \int_{a}^{b}|f(x)|dx \leqslant V(F, [a, b]) + \epsilon
\end{equation*}
Thus proving the other direction inequality.
\pagebreak

\section{The Stieltjes Integral}
We discussed the idea of the Riemann integral, in which we integrated with respect to partitions on the $x$-axis. Now, we integrate with respect to another function, we call the integrator:
\begin{defn}
A function $f: [a, b] \xrightarrow{} \mathbb{R}$ is Stieltjes integrable with respect to $g: [a, b] \xrightarrow{} \mathbb{R}$ if:
\begin{equation*}
  \int_{a}^{b}f dg \coloneq \lim_{||\Pi|| \xrightarrow{} 0}\sum_{i=1}^{n}f(t_i^*)(g(t_i) - g(t_{i-1})) \ \text{exists}
\end{equation*}
\end{defn}
Note that there are no restrictions (yet) on the behavior of $g$, so $g$ could have discontinuities, could be increasing, decreasing, or neither, etc. We will however observe one restriction later on. An equivalent way of restating this criterion is:
\begin{equation*}
  \forall \epsilon > 0, \exists \delta > 0, \forall \Pi: ||\Pi|| < \delta \implies \left|S(f, g, \Pi) - \int_{a}^{b}fdg\right| < \epsilon
\end{equation*}

If we write that $f \in RS(g, [a, b])$, it means that $f$ is Stieltjes integrable wrt $g$ on the interval $[a, b]$. Additionally, the notation $S(f, g, \Pi)$ is a RS-sum over the partition $\Pi$. We first note additivity properties of both $f$ and $g$:
\begin{lem}
Consider $f_1, f_2 \in RS(g, [a, b])$ and $\alpha, \beta \in \mathbb{R}$, then:
\begin{equation*}
  \int_{a}^{b}(\alpha f_1 + \beta f_2) dg = \alpha \int_{a}^{b}f_1 dg + \beta \int_{a}^{b}f_2 dg
\end{equation*}
\end{lem}
\begin{proof}
Because $f_1$ and $f_2$ are independently RS-integrable, we have the following:
\begin{align*}
&\forall \epsilon > 0, \exists \delta_1 > 0, \forall \Pi: ||\Pi|| < \delta_1 \implies \left|S(f_1, g, \Pi) - \int_{a}^{b}fdg\right| < \epsilon\\
&\forall \epsilon > 0, \exists \delta_2 > 0, \forall \Pi: ||\Pi|| < \delta_1 \implies \left|S(f_2, g, \Pi) - \int_{a}^{b}fdg\right| < \epsilon
\end{align*}
Then choosing $0 < \delta < \min\{\delta_1, \delta_2\}$, we get that $\forall \Pi: ||\Pi|| < \delta$:
\begin{equation*}
  S(\alpha f_1+\beta f_2, g, \Pi) = \sum_{i=1}^{n}(\alpha f_1(t_i^*) + \beta f_2(t_i^*))(g(t_i) - g(t_{i-1})) = \alpha \sum_{i=1}^{n}f_1(t_i^*)(g(t_i) - g(t_{i-1})) + \beta \sum_{i=1}^{n}f_2(t_i^*)(g(t_i) - g(t_{i-1}))
\end{equation*}
Then we get that:
\begin{equation*}
  \left|S(\alpha f_1 + \beta f_2, g, \Pi) - \alpha \int_{a}^{b}f+1dg - \beta \int_{a}^{b}f_2dg\right| = \left|\alpha (S(f_1, g, \Pi) - \int_{a}^{b}f_1dg) + \beta (S(f_2, g, \Pi) - \int_{a}^{b}f_2dg)\right|
\end{equation*}
\begin{equation*}
  \leqslant |\alpha| \cdot \epsilon + |\beta| \cdot \epsilon = \epsilon(|\alpha| + |\beta|)
\end{equation*}
Thus proving the criterion proposed in Definition 0.8 for RS-integrability.
\end{proof}
One additional lemma we did not encounter in terms of Riemann Integrability is the following:
\begin{lem}
Consider $f \in RS(g, [a, b]) \cap RS(h, [a, b])$, then:
\begin{equation*}
\int_{a}^{b}f d(\alpha g + \beta h) = \alpha \int_{a}^{b}f dg + \beta \int_{a}^{b}f dh
\end{equation*}
\end{lem}
\begin{proof}
We note that $f$ is RS-integrable with respect to both $g$ and $h$, each of which provides us $\delta_g, \delta_h$. Choosing $0 < \delta < \min\{\delta_g, \delta_h\}$, we see that:
\begin{equation*}
  S(f, \alpha g + \beta h, \Pi) = \sum_{i=1}^{n}f(t_i^*)(\alpha(g(t_i) - g(t_{i-1})) + \beta(h(t_i) - h(t_{i-1})))
\end{equation*}
\begin{equation*}
  \alpha \sum_{i=1}^{n}f(t_i^*)(g(t_i) - g(t_{i-1})) + \beta \sum_{i=1}^{n}f(t_i^*)(h(t_i) - h(t_{i-1}))
\end{equation*}
Once again applying the triangle inequality and RS-integrability wrt $g$ and $h$, we get the same bound as Lemma 0.14, proving the statement.
\end{proof}

Note an interesting consequence of this definition, which is that it does not necessarily tell us that $f$ is bounded. Consider a function unbounded within some interval. We can simply set $g$ to be constant in that interval, effectively canceling out its contributions to the RS-integral.\\

There is a case we will go over in a later case, perhaps involving measure. Another interesting lemma follows, which provides a restriction on the relationship between $f$, $g$. First we prove an oscillation characterization:
\begin{lem}
If $f \in RS(g, [a, b])$, then:
\begin{equation*}
  \forall \epsilon > 0, \exists \delta > 0, \forall \Pi: ||\Pi|| < \delta \implies \sum_{i=1}^{n}\text{osc}(f, [t_{i-1}, t_i])|g(t_i) - g(t_{i-1})| < \epsilon
\end{equation*}
\end{lem}
\begin{proof}
Let $\Pi, \Pi'$ be two marked partitions with mesh less than $\delta$ and the same subintervals. Then:
\begin{equation*}
  |S(f, g, \Pi) - S(f, g, \Pi')| = \sum_{i=1}^{n}|f(t_i^*) - f(t_i')|g(t_i) - g(t_{i-1})|
\end{equation*}
where we select the tagged points such that $\frac{1}{2}\text{osc}(f, [t_{i-1}, t_i]) \leqslant |f(t_i^*) - f(t_i')|$. Since $f$ is RS-integrable (represent for brevity the integral as $I$), we have that:
\begin{equation*}
  |S(f, g, \Pi) - S(f, g, \Pi')| = |S(f, g, \Pi) - I + I - S(f, g, \Pi')| < \epsilon/2
\end{equation*}
assuming we set the $\epsilon$-bound for RS-integrability as $\epsilon/4$. Then, because oscillation does not take into account marked partitions:
\begin{equation*}
  \forall \Pi: ||\Pi|| < \delta \implies \sum_{i=1}^{n}\text{osc}(f, [t_{i-1}, t_i])|g(t_i) - g(t_{i-1})| < \epsilon
\end{equation*}
We can further extend the ideas in this proof to derive a necessary and sufficient condition, a Cauchy Criterion of RS-integrability. First note however this is only a necessary condition, not a sufficient one.\\

This is due to the fact that proving RS-integrability requires marked partitions; however, such a condition only deals with unmarked partitions and oscillations. \textbf{THIS IS NOT A SATISFYING EXPLANATION, I WILL COME BACK TO THIS.}
\end{proof}
We now use Lemma 0.16 for the following statement:
\begin{lem}
If $f \in RS(g, [a, b])$, then $f$ and $g$ CANNOT share points of discontinuity.
\end{lem}
\begin{proof}
Assume that there exists some $x \in [a, b]$ such that both $f, g$ are discontinuous at $x$. Then:
\begin{align*}
  &\exists \epsilon > 0, \forall \delta > 0, \exists y \in [a, b]: |y - x| < \delta \land |f(y) - f(x)| \geqslant \epsilon\\
  &\exists \epsilon' > 0, \forall \delta > 0, \exists y' \in [a, b]: |y' - x| < \delta \land |g(y') - g(x)| \geqslant \epsilon'
\end{align*}
If we select $0 < \tilde{\epsilon} < \min\{\epsilon, \epsilon'\}$, then we can consolidate the two inequalities:
\begin{align*}
  &\exists \tilde{\epsilon}, \forall \delta > 0: \text{osc}(f, [x + \delta]) \geqslant \tilde{\epsilon} \lor \text{osc}(f, [x - \delta, x]) \geqslant \tilde{\epsilon}\\
  &\exists \tilde{\epsilon}, \forall \delta > 0, \exists y' \in [a, b]: |y' - x| < \delta \land |g(y') - g(x)| \geqslant \tilde{\epsilon}
\end{align*}
We know that a function is continuous if and only if it is left and right continuous. So in fact, the proof breaks into two cases:
\begin{enumerate}
\item $f, g$ are both NOT right (or left, WLOG) continuous at $x$
\item $f$ is NOT right continuous and $g$ is NOT left continuous (or vice versa)
\end{enumerate}
The first characterization for $f$ comes from the fact that if $f$ is not continuous, then:
\begin{equation*}
\tilde{\epsilon} \coloneq lim_{\delta \xrightarrow{} 0^+}\text{osc}(f, [x - \delta, x + \delta]) > 0
\end{equation*}
as $m_f(x) \neq M_f(x)$. This allows us to omit the idea of intermediate points for $f$, solely focusing on intervals.

\subsubsection*{Case One}

We also do not have monotonicity of $g$, however we do have monotonicity of the oscillation. Then we want to show an existence of a partition with mesh less than delta, where the oscillation sum is greater than some epsilon value.\\

Consider $n\delta > b-a$, then define $\Pi \coloneq \{[a + i \cdot \frac{b-a}{n}, a + (i+1) \cdot \frac{b-a}{n}] : 0 \leqslant i \leqslant n-1\}$. Then refine the partition to add $x$, the common point of discontinuity. Now take the partition point $t_j$ closest to $x$ and set $\delta = |t_j - x|$. Then within this $\delta$, we have the point $y'$ for the discontinuity of $g$.\\

We refine the parition further, adding $y'$ to the partition points, getting an interval $[\min\{y', x\}, \max\{y', x\}]$. Now we know that $f, g$ are BOTH not right continuous at $x$, so setting $\delta = |y' - x|$:
\begin{equation*}
  \text{osc}(f, [x, y']) \geqslant \tilde{\epsilon}
\end{equation*}
So, we get that:
\begin{equation*}
  \forall \Pi: ||\Pi|| < \delta \implies \sum_{i=1}^{n}\text{osc}(f, [t_{i-1}, t_i])|g(t_i) - g(t_{i-1})| \geqslant \tilde{\epsilon}^2
\end{equation*}
So the discontinuities cannot appear on the same side of the functions. Now what about case 2, based on this proof, seems possible right?

\subsubsection*{Case Two}
Now instead of adding the points of discontinuity to the partition themselves, we will think in terms of partition subintervals. This proof would be easier if we had monotonicity of $g$, as considering some subtinterval such that $x \in [t_j, t_{j+1}]$, we would have some $y'$ in the interval. Using monotonicity of oscillation and the fact that $x$ is an interior point, we could say:
\begin{equation*}
  \exists \delta > 0: [x - \delta, x + \delta] \subseteq [t_j, t_{j+1}] \implies \text{osc}(f, [t_j, t_{j+1}]) \geqslant \text{osc}(f, [x - \delta, x + \delta]) > \tilde{\epsilon}
\end{equation*}
Then we simply have $|g(t_{j+1}) - g(t_j)| \geqslant |g(y') - g(x)|$ because of monotonicity, so the oscillation sum fails to be less than all epsilon. In this case, we did not need to make a distinction between left/right continuous because it did not matter which side the $y'$ appeared on (we did not have to refine any partition).\\

Now assume there exists some subinterval such that $x \in [t_j, t_{j+1}]$ (this must be true by covering properties). Then we can guarantee some $y' \in [t_j, x]$ by non-left continuity of $g$. Now assume $f \in RS(g, [a, b])$, with a partition $\Pi$ such that $||\Pi|| < \delta$. Then:
\begin{equation*}
  \epsilon > \sum_{i=1}^{n}\text{osc}(f, [t_{i-1}, t_i])|g(t_i) - g(t_{i-1})| \geqslant \text{osc}(f, [t_j, t_{j+1}])|g(t_{j+1}) - g(t_j)| \geqslant \tilde{\epsilon}|g(t_{j+1}) - g(t_j)|
\end{equation*}
Because $f$ is not right continuous, this is true of any $t_j \coloneq x - \delta', t_{j+1} \coloneq x + \delta'$ where $0 < 2\delta' < \delta$. So we get that:
\begin{equation*}
  \forall \delta' > 0: 0 < |t_{j+1} - t_j| \leqslant 2\delta' \implies |g(t_{j+1}) - g(t_j)| < \epsilon/\tilde{\epsilon} \implies \text{osc}(g, [t_j, t_{j+1}]) \leqslant \epsilon/\tilde{\epsilon}
\end{equation*} 
We can make this construction with $\delta'$ as integrability must hold for ANY $\Pi$ as long as its mesh $||\Pi|| < \delta$. So, the oscillation being bounded means:
\begin{equation*}
  \forall \epsilon > 0, \exists \delta' > 0: |s-t| \leqslant 2\delta' \implies \text{osc}(g, [s, t]) < \epsilon/\tilde{\epsilon}
\end{equation*}
So $m_g(x) = M_g(x)$, thus $g$ would be continuous at $x$, which contradicts our initial condition.
\end{proof}
NOTE: An interesting consequence of this proof is that if $f \in RS(g, [a, b])$, then $f \in RS(g, [a, c])$ and $f \in RS(g, [c, d])$. However the converse is not necessarily true. It is possible for $f \in RS(g, [a, c])$, with $f$ being NOT right continuous at $c$ and $g$ being NOT left continuous.\\

On each individual interval, they do not share a discontinuity; however on all of $[a, b]$, we obtain Case 2 in Lemma 0.17, which is not possible. We now provide a sufficient and necessary condition, the Cauchy Criterion for RS-integrability:
\begin{lem}
$f \in RS(g, [a, b])$ if and only if:
\begin{equation}
\forall \epsilon > 0, \exists \delta > 0, \forall \Pi: \max\{||\Pi||, ||\Pi'||\} < \delta \implies |S(f, g, \Pi) - S(f, g, \Pi')| < \epsilon
\end{equation}
\end{lem}
\begin{proof}
The first part of the proof is quite straightforward, similar to what we used in Lemma 0.16 (it comes from triangle inequality). The converse follws from the completeness of $\mathbb{R}$. While we can make a convergence of nets argument, we choose not to here. Let $\Pi, \Pi'$ be partitions with the same subintervals, but different tags.
\end{proof}
This Lemma actually provides a good explanation as to why the converse of Lemma 0.16 is not necessarily true. Lemma 0.16 relies on one unmarked partition only, so the maximum freedom we get in the converse is choosing the marked points. However by this Cauchy Criterion, it must be true for ANY two intervals with max mesh less than $\delta$ (which means even if their subintervals are not aligned). So the result of the converse of Lemma 0.16 is too weak to prove integrability.

\pagebreak
\section{Pointwise and Uniform Convergence}
We now consider sequences of functions $\{f_n\}_{n \in \mathbb{N}}$ and want to see if and how they converge to some function $f$. The important of these falls under measure theory and integration theory:
\begin{defn}
Given two metric spaces $(X, \varrho)$ and $(Y, \varrho)$, a sequence of functions $\{f_n\}_{n \in \mathbb{N}}$ that map from $X$ to $Y$ is pointwise convergent if:
\begin{equation*}
\forall \epsilon > 0, \forall x \in X, \exists n_0 \geqslant 0, \forall n \geqslant n_0: \varrho(f_n(x), f(x)) < \epsilon
\end{equation*}
\end{defn}
Now akin to the idea of uniform continuity, we introduce the idea of uniform convergence:
\begin{defn}
A sequence $\{f_n\}_{n \in \mathbb{N}}$ is called uniformly convergent if:
\begin{equation*}
\forall \epsilon > 0, \exists n_0 \geqslant 0, \forall x \in X, \forall n \geqslant n_0: \varrho(f_n(x), f(x)) < \epsilon
\end{equation*}
\end{defn}
Note the difference here is the quantifier switch. One $n_0$ (similar to one $\delta$) works for all possible $x$ in the domain, yielding a stronger condition. As it is for continuity, we view the following lemma:
\begin{lem}
A uniformly convergent sequence $\{f_n\}_{n \in \mathbb{N}}$ is also pointwise convergent.
\end{lem}
\begin{proof}
The proof follows from a simple quantifier switch.
\end{proof}

\medskip
\subsection{Examples}
We now provide some examples of function sequences that are pointwise and uniformly convergent:\\

\textbf{Pointwise}: Consider the sequence, where:
\begin{equation*}
  f_n(x) \coloneq \frac{n}{1+nx} \land f_n: (0, 1) \xrightarrow{} \mathbb{R}
\end{equation*}
We claim that this sequence converges to $f(x) = 1/x$, and is pointwise convergent. So $\forall \epsilon > 0$:
\begin{equation*}
\left|\frac{n}{1+nx} - \frac{1}{x}\right| = \left|\frac{1}{x(1+nx)}\right| \leqslant \left|\frac{1}{nx}\right| < \epsilon
\end{equation*}
So as long as we have that $n > 1/(\epsilon \cdot x)$, we know that the further functions in the sequence are within an epsilon from $1/x$. Notice how in our representation for $n$, we have $x$, showing it is dependent on what value in the domain we are looking at.\\

One way we can show is by the $\epsilon$-definition; however, another way we can show it is the limsup definition or the negation of it. We want to show:
\begin{equation*}
  \exists \epsilon > 0, \forall n \in \mathbb{N}: sup_{x \in (0, 1)}|f_n(x) - \frac{1}{x}| \geqslant \epsilon
\end{equation*}
Then we choose $\epsilon = 1$. We see that:
\begin{equation*}
  \forall n \in \mathbb{N}: \frac{1}{1+nx} \leqslant \frac{1}{x(1+nx)} \implies \sup_{x \in (0, 1)}\frac{1}{1+nx} \leqslant \sup_{x \in (0, 1)}\frac{1}{x(1+nx)}
\end{equation*}
The supremum of the smaller term is always $1$ for all $n \in \mathbb{N}$, so we have proven our statement.\\

The interesting thing about this example is that $\forall n \in \mathbb{N}$, we have that $f_n$ are bounded, yet the limit of the function $1/x$ is unbounded:
\begin{equation*}
  \forall x \in (0, 1): 1 + nx > 1 \implies \left|\frac{n}{1+nx}\right| < n
\end{equation*}
What this tells us is that pointwise convergence does not preserve bounding properties.\\

We look at another example, where pointwise convergence does not preserve continuity either. Consider the sequence:
\begin{equation*}
  f_n(x) \coloneq x^n \land f: [0, 1] \xrightarrow{} \mathbb{R}
\end{equation*}
We prove this is a pointwise convergent sequence, which we claim converges to $0$ for all $x \in [0, 1)$. So $\forall \epsilon > 0$:
\begin{equation*}
  |x^n| < \epsilon \implies n > \frac{\ln(\epsilon)}{\ln(x)}
\end{equation*}
Once again it depends on $x$, showing it is not uniformly convergent. Now the interesting behavior we want to analyze here is that each function $f_n$ is continuous due to it being a polynomial. However, the limit which is $0$ everywhere in $[0, 1)$ but is $1$ at $x = 1$ is not left continuous at $x = 1$. So continuity is not preserved by pointwise convergence.\\

Now we examined two sequences which are both pointwise convergent, one being bounded but the boundedness property not transferring, and the other being continuous, but continuity not transferring. So, we view the following lemma:
\begin{lem}
If $\forall n \in \mathbb{N}, \exists M_n, \forall x \in X: |f_n(x)| \leqslant M_n$ and $f_n \xrightarrow{} f$ uniformly, then $f(x)$ is bounded.
\end{lem}
\begin{proof}
As we do many of these proofs we will notice that they utilize the triangle inequality. Now, for say $\epsilon = 1$, uniform continuity gives us an $n_0$ that satisfies the criterion on all of $X$. Then:
\begin{equation*}
  |f(x)| = |f(x) - f_{n_0}(x) + f_{n_0}(x)| \leqslant |f(x) - f_{n_0}(x)| + |f_{n_0}(x)| \leqslant 1 + M_{n_0}
\end{equation*}
Thus $f$ is bounded. Now if we did not have uniform continuity, then the same $n_0$ may not work at some other $x' \in X$. In which case, our bound at $x'$ might be something different like $1 + M_{n_1}$ which could be greater then the previous, so it may be possible that $f$ is greater than $1 + M_{n_0}$.
\end{proof}

Now we introduce a stronger theorem that talks about the preservation of continuity:
\begin{thm}
If $f_n \xrightarrow{} f$ uniformly, then if all $f_n$ are continuous at $x_0$, $f$ is continuous at $x_0$.
\end{thm}
\begin{proof}
We use the epsilon-delta definition of continuity:
\begin{equation*}
  \forall n \in \mathbb{N}, \forall x_0 \in X, \forall \epsilon > 0, \exists \delta > 0, \forall x \in X: |x-x_0| < \delta \implies |f_n(x) - f_n(x_0)| < \epsilon
\end{equation*}
\begin{equation*}
  \forall \epsilon > 0, \exists n_0 \geqslant 0, \forall x \in X, \forall n \geqslant n_0: |f_n(x) - f(x)| < \epsilon
\end{equation*}
Now, we write the following via triangle inequality expansion:
\begin{equation*}
\varrho(f(x), f(x_0)) \leqslant \varrho(f(x), f_n(x)) + \varrho(f_n(x), f_n(x_0)) + \varrho(f_n(x_0), f(x_0)) < 3\epsilon
\end{equation*}
Taking $n_0$ from uniform convergence, this inequality works for all $n \geqslant n_0$ as long as we use its $\delta_n$ provided by continuity (because each $n$ produces its own $\delta$ for some $\epsilon$).\\

Note that pointwise convergence would not be enough because we would have that for each $x \in X$, there exists some $\delta_x$. Instead, we want it so that there exists a $\delta$ such that for each $x \in X$ which is delta away from $x_0$ at most... follows from definition of continuity.
\end{proof}
We can replace the hypothesis of continuity in Theorem 0.3 with that of the limit existing, but proving this requires the extra assumption of the completeness of $Y$ (this is what allows us to interchange limits). We can also weaken the idea of convergence, to bring about the following two ideas:
\begin{defn}
A sequence of functions is called pointwise Cauchy if:
\begin{equation*}
\forall x \in X, \forall \epsilon > 0, \exists n_0 \geqslant 0, \forall m, n \geqslant n_0: \varrho(f_n(x), f_m(x)) < \epsilon
\end{equation*}
Similarly such a sequence is called uniformly Cauchy if:
\begin{equation*}
\forall \epsilon > 0, \exists n_0 \geqslant 0, \forall m, n \geqslant n_0, \forall x \in X: \varrho(f_n(x), f_m(x)) < \epsilon
\end{equation*}
\end{defn}
We now view a lemma that will yield us an interesting result (allowing us to weaken sufficient conditions):
\begin{lem}
If $f_n \xrightarrow{} f$ converges uniformly, then $f_n$ is uniformly Cauchy. For the converse, if $f_n$ is uniformly Cauchy, then if $Y$ is complete, $f_n$ converges uniformly. 
\end{lem}
\begin{proof}
$\implies$: While not an entirely symmetric proof, we can add completeness in this direction without changing any result; but we choose to weaken it. If $f_n$ is uniformly convergent, then:
\begin{equation*}
  \forall \epsilon > 0, \exists n_0 \geqslant 0, \forall n \geqslant n_0, \forall x \in X: \varrho(f_n(x), f(x)) < \epsilon
\end{equation*}
Then $\forall m, n \geqslant n_0$:
\begin{equation*}
  \varrho(f_m(x), f_n(x)) \leqslant \varrho(f_m(x), f(x)) + \varrho(f_n(x), f(x)) < 2\epsilon
\end{equation*}
Thus proving it to be Cauchy. Now for the converse:
$\impliedby$: If $f_n$ is uniformly Cauchy, we have that it is pointwise Cauchy (this comes from a simple quantifier swap), which in a complete space means its pointwise convergent to some $f(x)$. Then:
\begin{equation*}
\forall \epsilon > 0, \exists n_0 \geqslant 0, \forall n \geqslant n_0, \forall x \in X: \varrho(f_n(x), f(x)) \leqslant \varrho(f_n(x), f_{n_0}(x)) + \varrho(f_{n_0}(x), f(x)) < 2\epsilon
\end{equation*}
Note that here we only needed the condition of pointwise convergence, which is a weaker sufficient condition than completeness of $Y$. 
\end{proof}

\subsection{Equicontinuity}
Equicontinuity is the property that allows for all functions $f_n$ in a sequence to be continuous "at the same rate" which means they have the same $\delta$ of continuity. Formally:
\begin{equation*}
  \forall x_0 \in X, \forall \epsilon > 0, \exists \delta > 0, \forall n \in \mathbb{N}: \varrho(x, x_0) < \delta \implies \varrho(f_n(x), f_n(x_0)) < \epsilon
\end{equation*}
We can extend this idea to uniform continuity to get uniform equicontinuity, which follows from a simple switch of the quantifiers. Consider the following family:
\begin{equation*}
  \forall n \in \mathbb{N}: f_n(x) \coloneq x^2 + \frac{(-1)^n}{n+1}
\end{equation*}
We first prove, by the supremum characterization, that this sequence is uniformly convergent:
\begin{equation*}
  \sup\left|x^2 + \frac{(-1)^n}{n+1} - x^2\right| = \sup\left|\frac{1}{n+1}\right| \xrightarrow{} 0
\end{equation*}
All members of this family is convex because their second derivative is greater than or equal to $0$. Now we know that the derivative of the second lines is increasing across the domain of a convex function; however, we consider a compact subinterval of the domain $[b, c]$ to get uniform equicontinuity.\\

We previously derived the following property of convex functions. Consider two points $a, b$ such that $a < b < c < d$, then:
\begin{equation*}
\forall x < y \in [b, c]: \frac{f(a) - f(b)}{a-b} < \frac{f(b) - f(x)}{b-x} < \frac{f(x) - f(y)}{x-y} < \frac{f(y) - f(c)}{y-c} < \frac{f(c) - f(d)}{c-d}
\end{equation*}
Now for locally lipschitz continuity, we require some constant $C$. Notice how $(f(x) - f(y))/(x-y)$ is stuck between the fractions consisting of $a, b$ and $c, d$, which represent the slopes of the secant lines. Now, it is possible that the leftmost slope is negative, but greater in magnitude than the rightmost, so if we choose:
\begin{equation*}
  C \coloneq \max\left\{\left|\frac{f(a) - f(b)}{a-b}\right|, \left|\frac{f(c) - f(d)}{c-d}\right|\right\}
\end{equation*}
Then we see that:
\begin{equation*}
  \forall x, y \in [b, c]: \left|\frac{f(x) - f(y)}{x-y}\right| < C \implies |f(x) - f(y)| < C|x-y|
\end{equation*}
Thus proving our local lipschitz condition as long as we set $\delta = \epsilon/C$ ONLY for points within the interval $[b, c]$. Now examining our family $f_n$, we see that the $n$-term cancels out, so the same $C$ works for all $f_n$, thus the same $\delta$ works, showing it is a uniformly equicontinuous family.\\

Notice the pattern we saw in the previous example. Each $f_n$ was uniformly continuous, and with the sequence being uniformly convergent, we see that they are uniformly equicontinuous. Thus:
\begin{lem}
$f_n \xrightarrow{} f$ uniformly and each $f_n$ is continuous, then $f_n$ is equicontinuous. 
\end{lem}
\begin{proof}
First by Theorem 0.3, we see that $f$ must be continuous, which means:
\begin{equation*}
  \forall \epsilon > 0, \exists \delta' > 0: \varrho(x, x_0) < \delta_0 \implies \varrho(f(x), f(x_0)) < \epsilon
\end{equation*}
Additionally by uniform convergence:
\begin{equation*}
  \forall \epsilon > 0, \exists n_0 \geqslant 0, \forall n \geqslant n_0: \varrho(f_n(x), f(x)) < \epsilon
\end{equation*}
Then we see that $\forall x \in B(x, \delta')$:
\begin{equation*}
  \varrho(f_n(x), f_n(x_0)) \leqslant \varrho(f_n(x), f(x)) + \varrho(f(x), f(x_0))  + \varrho(f(x_0), f_n(x_0)) < 3\epsilon
\end{equation*}
Thus we have shown that this $\delta_0$-ball of continuity extends to $f_n$ for $n \geqslant n_0$. Now, for $n < n_0$, we choose:
\begin{equation*}
  \delta'' \coloneq \min{\delta_n: 0 \leqslant n < n_0}
\end{equation*}
Then if we choose $\delta = \min\{\delta', \delta''\}$, we find that we have equicontinuity for all $n \in \mathbb{N}$. Now notice how the idea of uniform convergence pulls weight in the proof.\\

Without uniform convergence, the $n_0$ for convergence would be different for each $x \in X$, which first would mean we could not claim continuity of $f$ AND we cannot claim equicontinuity of $f_n$ because we would have to take different minimums for each $x$.
\end{proof}

We can now extend this to uniform equicontinuity with the following corollary:
\begin{cor}
$f_n \xrightarrow{} f$ uniformly and each $f_n$ uniformly continuous, then $f_n$ is uniformly equicontinuous.
\end{cor}
\begin{proof}
Now uniform convergence gives us the stronger condition that:
\begin{equation*}
  \forall n \in \mathbb{N}, \forall \epsilon > 0, \exists \delta_n > 0, \forall x, y \in X: \varrho(x, y) < \delta_n \implies \varrho(f_n(x), f_n(y)) < \epsilon
\end{equation*}
Now we first need to prove that $f$ is uniformly continuous, with $n_0$ coming from uniform convergence:
\begin{equation*}
  \varrho(f(x), f(y)) \leqslant \varrho(f(x), f_{n_0}(x)) + \varrho(f_{n_0}(x), f_{n_0}(y)) + \varrho(f_{n_0}(y), f(y))
\end{equation*}
Using $\delta_{n_0}$, we see that $\varrho(f(x), f(y)) < 3\epsilon$, thus we can extend Theorem 0.3 to uniform continuity. Now to prove the Corollary:
\begin{equation*}
\forall \epsilon > 0, \exists \delta' > 0: \varrho(x, y) < \delta' \implies \varrho(f(x), f(y)) < \epsilon
\end{equation*}
From uniform convergence, we get that for $n \geqslant n_0$:
\begin{equation*}
  \varrho(f_n(x), f_n(y)) \leqslant \varrho(f_n(x), f(x)) + \varrho(f(x), f(y)) + \varrho(f(y), f_n(y)) < 3\epsilon
\end{equation*}
Using a similar minimum strategy for $n < n_0$, we have proven uniform equicontinuity.
\end{proof}

Now notice one thing in Lemma 0.22 and Corollary 0.3 here. We use $f$ as a bridge to prove equicontinuity by using its (uniform or pointwise) continuity. However, we already know that all of the $f_n$ share the same property, thus it is adequate to have the uniformly Cauchy property in which we simply use $f_{n_0}$ as a bridge.

\subsection{Families of Equicontinuous Functions}
The most common way to prove a family of functions is equicontinuous is to show that they admit the same modulus of continuity, which recall is a function $\omega: [0, \infty) \xrightarrow{} [0, \infty)$ such that:
\begin{enumerate}
  \item $\omega$ is increasing
  \item $|f(x) - f(y)| \leqslant \omega(|x-y|) \iff f \ \text{uniformly continuous}$
  \item $\lim_{t \xrightarrow{} 0^+} \omega(t) = 0$
\end{enumerate}
In the case of Lipschitz continuous (or H\"older continuous) functions, we aim to show there is some common constant $C$ that acts as the modulus of continuity. We can show this by a uniform bound on the derivative.\\

For example the following family of polynomials on $[a, b]$ up to the $n$-th degree is equicontinuous:
\begin{equation*}
  \left\{x \mapsto \sum_{k=0}^{n}a_kx^k: a_0 \in \mathbb{R} \land a_1,...,a_k \in [-M, M]\right\}
\end{equation*}
We can prove this by analyzing the derivatives of each of the polynomial expressions:
\begin{equation*}
  \left|\frac{d}{dx}\sum_{k=0}^{n}a_kx^k\right| = \left|\sum_{k=1}^{n}k \cdot a_k x^{k-1}\right| \leqslant \left|M\sum_{k=1}^{n}k \cdot x^{k-1}\right| \leqslant \left|M\sum_{k=1}^{n}k \cdot b^{k-1}\right|
\end{equation*}
Therefore, we have a bound on the derivative for all members in this family, we will call $\lambda$ represented by the expression on the right. Now consider any $x, y \in [a, b]$, specifically the interval $[x, y]$. For any $f$ within the family, we know that it is continuous on $[a, b]$ and differentiable on $(a, b)$ because it is a polynomial.\\

Thus, we can apply MVT on $[x, y]$ and say that:
\begin{equation*}
  \exists c \in [x, y]: f'(c) = \frac{f(y) - f(x)}{y-x} \implies |f(y) - f(x)| = |f'(c)||y-x| \leqslant \lambda|y-x|
\end{equation*}
Since $f$ and $x, y$ were arbitrary, we have uniform continuity for each $f$ in the family. Furthermore, because the same $\lambda$ applies to all such polynomials in the family, we have uniform equicontinuity.\\

This idea can be extended to show the following family in greater generality is also uniformly equicontinuous:
\begin{equation*}
\left\{f \in \mathbb{R}^{[a, b]}: f \ \text{continuous on } [a, b] \land f \ \text{differentiable on } (a, b) \land |f'(x)| < M\right\}
\end{equation*}

Now consider another example of the following sequence with $f_n: \mathbb{Q} \xrightarrow{} \mathbb{R}$:
\begin{equation*}
  f_n(x) \coloneq \begin{cases*}
    1 \ : \ nx < \sqrt{2}\\
    0 \ : \ nx > \sqrt{2}
  \end{cases*}
\end{equation*}
We show that $f_n(x)$ is not uniformly convergent NOR uniformly Cauchy. Each $f_n$ is however continuous and the family itself is equicontinuous, so the sufficient conditions of weak Lemma 7.4 with uniform Cauchy are not satisfied. We claim it converges to $0$-function.\\

First we prove it is not uniformly convergent by the supremum categorization:
\begin{equation*}
  \forall n_0 \geqslant 0, \exists n = n_0, \exists x \in \mathbb{Q}: 0 < x < \frac{\sqrt{2}}{n} \implies \sup_{x \in \mathbb{Q}}|f_n(x) - 0| = 1 \nrightarrow 0
\end{equation*}
which follows from density of the rationals. We now prove it is not uniformly Cauchy:
\begin{equation*}
  \exists \epsilon > 0, \forall n_0 \geqslant 0, \exists m, n \geqslant n_0, \exists x \in \mathbb{Q}: |f_n(x) - f_m(x)| \geqslant \epsilon
\end{equation*}
If we set $\epsilon = 1/2$, then for any $n_0 \geqslant 0$, we take $m > n = n_0$, then:
\begin{equation*}
  \exists x \in \mathbb{Q}: \frac{\sqrt{2}}{m} < x < \frac{\sqrt{2}}{n} \implies |f_n(x) - f_m(x)| = |0 - 1| = 1 > 1/2
\end{equation*}
As a side effort, we also prove the continuity of each $f_n$, which follows from the $\epsilon-\delta$ definition:
\begin{equation*}
  \forall \epsilon > 0, \exists \delta > 0: |x - x_0| < \delta \implies |f_n(x) - f_n(x_0)| < \epsilon
\end{equation*}
We now have a jump, but because our domain is $\mathbb{Q}$, the function still remains continuous. If we have $x_0 < \sqrt{2}/n$, then we choose $\delta = |\sqrt{2}/n - x|/2$. The same goes for the case when $x_0 > \sqrt{2}/n$. In these $\delta$-balls, the function values agree.\\

Finally, we prove the family is equicontinuous for the restricted domain $\mathbb{Q} \cap (0, 1)$:
\begin{equation*}
  \forall \epsilon > 0, \exists \delta > 0: |x-x_0| < \delta \implies |f_n(x) - f_n(x_0)| < \epsilon
\end{equation*}
If we choose $\delta \coloneq \min_{n > 2}|x_0 - \sqrt{2}/n|/2$, then we are guaranteed that we end up on one side of the closest jump; so for any further jumps, we are guaranteed that we remain ending up on only one side of the cut (because $\delta$ is a minimum distance).\\

Now consider the same family of functions, but on the domain $\mathbb{Q} \cap (1, 2)$ v.s. the domain $\mathbb{Q} \cap [1, 2]$. We once again in a similar fashion can prove that $f_n$ is equicontinuous in this interval and that it converges pointwise; however, in the previous example, the point $0$ is a limit point of a sequence of shape $1/n$, so we would lose this equicontinuous structure if we considered the domain $\mathbb{Q} \cap [0, 1]$.\\

However, neither $1$ nor $2$ are limit points. Now, we can prove in the same fashion that the sequence is. (this does not work) 

\subsection{Mini Arzela-Ascoli}
\begin{thm}
Given $f_n: X \xrightarrow{} Y$, we have that:
\begin{equation*}
  \{f_n\}_{n \in \mathbb{N}} \ \text{equicontinuous} \land f_n \xrightarrow{} f \ \text{pointwise} \land X \ \text{compact}
\end{equation*}
\begin{equation*}
  \implies f_n \xrightarrow{} f \ \text{uniformly}
\end{equation*}
\end{thm}
\begin{proof}
Now given $f_n$ is equicontinuous, we know that:
\begin{equation*}
\forall x_0 \in X, \forall \epsilon > 0, \exists \delta > 0, \forall n \in \mathbb{N}: \varrho(x , x_0) < \delta \implies \varrho(f_n(x), f_n(x_0)) < \epsilon
\end{equation*}
We then define:
\begin{equation*}
  \forall x \in X: \delta(x) \coloneq \frac{1}{2}\sup\{\delta \in (0, 1): (\forall n \in \mathbb{N}: f_n(B(x, \delta)) \subseteq B(f(x), \epsilon))\}
\end{equation*}
Because of equicontinuity, we have that $\delta(x)$ is always nonempty (because it is a subset of a delta-ball provided by equicontinuity). Now the set of $B(x, \delta(x))$ acts as an open cover for $X$.\\

Since $X$ is compact, it admits a finite subcover with points $x_0,...,x_m$. Now, we can apply a $3\epsilon$ proof. Pointwise continuity gives us an $n_i$ at each $x_i$, so choosing $n' = \max\{n_i : i \in 0,...,m\}$:
\begin{equation*}
  \forall m, n \geqslant n', \forall x \in X, \exists i \in 0,...,m: x \in B(x_i, \delta(x_i)) \land \varrho(f_n(x), f_m(x)) \leqslant \varrho(f_n(x), f_n(x_i)) + \varrho(f_n(x_i), f_m(x_i)) + \varrho(f_m(x_i)) < 3\epsilon
\end{equation*}
where the middle term is taken care by pointwise convergence implying pointwise Cauchy. Having shown uniformly Cauchy, in a complete (or compact) space, this can be upgraded to uniform convergence.
\end{proof}

We can switch the sufficient conditions of this theorem to get the following:
\begin{cor}
Given $f_n: X \xrightarrow{} Y$, we have that:
\begin{equation*}
  \{f_n\}_{n \in \mathbb{N}} \ \text{ uniformly equicontinuous} \land f_n \xrightarrow{} f \ \text{pointwise} \land X \ \text{totally bounded}
\end{equation*}
\begin{equation*}
  \implies f_n \xrightarrow{} f \ \text{uniformly}
\end{equation*}
\end{cor}
\begin{proof}
Totally bounded gives us a finite cover. With the same $\delta(x)$-ball construction, we can do another $3\epsilon$ proof:
\end{proof}

\pagebreak
\section*{Metric Space of Functions}
While we have been using these notions in the 

\end{document}